% Generated by tools/textbook/export_listings.js from tools/concepts/assessments.json: do not edit.
\begin{chapstatus}{The thread of this chapter}
The machine you programmed in Lecture 02 is a von Neumann machine: program and data in one memory. RISC-V is the latest step of a long debate on how simple instructions should be. Next: how anything at all becomes bits.
\par\medskip\textbf{Exit ticket} (two to four minutes, alone, on paper):
\begin{enumerate}[nosep,leftmargin=2em]
\item Explain the sentence attributed to Dijkstra, 'Computer science is no more about computers than astronomy is about telescopes', using the definition of informatics given in the lecture.
\item Where is the program of a cartridge kept while it runs, and why does this make a von Neumann machine different from a machine programmed by rewiring cables?
\end{enumerate}
\end{chapstatus}
